\begin{tcolorbox}[enhanced,attach boxed title to top center=
	{yshift=-2mm,yshifttext=-1mm},colback=white,colframe=black,colbacktitle=black,title={\textbf{Instructions for the exam}}]
	\begin{itemize}
		\item You will have \textbf{\examtime\ minutes} to complete the exam.
		\item There are \textbf{\examproblemcount\ problems}, some of which may have multiple parts. Each problem is worth a total of \textbf{\examproblemweight\ points}.
		
		\item A graphing calculator is \emph{\graphingcalculator} for this exam.
		\item \textbf{This is a closed-book exam.} You may not refer to external sources such as notes (including the course notes) and books, and you are not allowed to receive help from anyone else.
		
		However, you \emph{are} allowed to use a notecard for the exam. For the notecard, you can only use \textbf{at most one sheet} with maximum dimensions of \textbf{5 inches by 8 inches}. You may use both sides of the notecard to fill as much material you'll reference during the exam.
		\item \textbf{No work, no credit!} Unless otherwise indicated, every problem requires work to be shown; that is, you'll need to give a valid explanation or process along with your answer. Such problems that include correct answers with no valid explanation will not be given any credit. Incorrect or incomplete answers where some correct work is shown are likely to receive some partial credit.
		\item \textbf{Need space to show work?} Each problem page has remaining space for your work and solution; this includes the empty back page of the problem page. If there is work that you do not want to be counted for grading, cross it out with a large ``X." If you use an additional page for any problem, please write a note on the problem page (example: ``Go to page 4" or ``Go to scrap paper") telling the grader where to look, and label your work on those pages for the designated problem as necessary (example: to the left of your work for Problem 4(a), write ``4(a)").
		
		You may ask for extra scrap paper as provided by the Testing Center or Learning Support Services, and whatever scrap paper you have used will be collected and attached to the exam. Your work on scrap paper will be not counted for grading, \emph{unless you indicate (on the original problem page) the grader to look over there, in order for your work to be counted}. Thus, even if you have a complete and correct explanation for your solution on scrap paper (and not anywhere else on the exam paper), you will \textbf{not} receive any credit if you don't redirect the grader to the respective scrap paper page(s) for that explanation.
		% Additional page policy from Testing Center?
	\end{itemize}
	\centering \Large{Good luck!} 
	\vspace*{1em}
\end{tcolorbox}